\section{Related Work}
\label{sec:related}

\subsection{EP Communication}
DeepEP~\cite{deepep} and NCCL EP~\cite{ncclep} provide GPU-initiated communication paths for expert parallelism. Their metadata exchanges, preallocated buffers, and synchronization schedules determine when data can be sent. \sys{} uses network staging to begin upload while peer readiness is still being collected.

SwiftEP~\cite{swiftep} optimizes buffer movement and GPU work, while COMET~\cite{comet} overlaps communication with computation. Perseus~\cite{perseus} reorganizes PUT, fence, and signal processing in the transfer pipeline. UBEP~\cite{ubep} uses point-to-point signals and Data-as-Flag to reduce stage synchronization in scale-up EP. These approaches improve data movement or its coordination; early upload complements them when a readiness dependency precedes transmission.

\subsection{In-Network Computation}
SHARP~\cite{sharp} and SwitchML~\cite{switchml} use network computation for collective reduction. DySHARP~\cite{dysharp} combines dynamic addressing, multicast, and reduction for MoE, while MultiWrite~\cite{multiwrite} uses multicast writes to reduce repeated upstream transfers. These systems demonstrate network processing and state management for distributed workloads. \sys{} applies temporary state and buffering to overlap pre-barrier coordination with upload.

\subsection{Synchronization and Delay Analysis}
EPIC~\cite{epic} provides protocol and resource abstractions for Ethernet in-network computing. GPU-initiated networking~\cite{gin} provides endpoint communication primitives. \sys{} combines these capabilities with readiness checks, symmetric receive buffers, and flow control.

Swift~\cite{swift} separates processing and network contributions to delay. We apply that distinction to the execution order of an EP operation. FabricPerf~\cite{fabricperf} further motivates explicit timing boundaries; our pre-barrier measurements pair each interval with the corresponding operator or GPU step.
